\en{In this chapter we will see that two lattices that are rotated and/or scaled versions of each other can be called "equivalent" and that equivalent lattices have the same value of Klein's absolute invariant~$J$.}%
\de{In diesem Kapitel werden wir sehen, dass zwei Gitter, die durch eine Drehstreckung auseinander hervorgehen, \glqq äqui"-valent\grqq~genannt werden können und dass äquivalente Gitter die gleiche absolute Invariante $J$ haben.}

\begin{defi}
\en{Two lattices $L,L'\subset \mathbb C$ are called "equivalent", iff they can be obtained from each other by rotation and scaling, i.e.~iff there is $a\in\mathbb C$ with $L'=a\cdot L$ and $a\neq 0$.}%
\de{Zwei Gitter $L$ und $L'$, die durch eine Drehstreckung $L'=a\cdot L$ mit $a\in\mathbb C$ auseinander hervorgehen $(a\neq 0)$, heißen äquivalent.}
\end{defi}

\begin{bem}
\en{Any elliptic function $f(z)$ of the lattice $L$ yields an elliptic function $g(z)=f\lk\frac z a\rk$ of the lattice $L'=a\cdot L$ and vice versa. That's why we call $L$ and $L'$ equivalent.}%
\de{Für jede elliptische Funktion $f(z)$ zum Gitter $L$ ist $g(z)=f\lk\frac z a\rk$ eine elliptische Funktion zum Gitter $L'=a\cdot L$. Deshalb nennt man $L$ und $L'$ auch äquivalent.}
\end{bem}

\begin{thm}\label{\en{EN}ltau}
\en{For each lattice $L=\mathbb Z \omega_1 + \mathbb Z \omega_2$ there is an equivalent lattice $L_\tau=\mathbb Z + \mathbb Z \tau$ with $\tau$ from the upper half plane $\mathbb H$ (i.e.~$\im(\tau)>0$).}%
\de{Zu jedem Gitter $L=\mathbb Z \omega_1 + \mathbb Z \omega_2$ gibt es ein äquivalentes Gitter $L_\tau=\mathbb Z + \mathbb Z \tau$, wobei $\tau$ in der oberen Halbebene $\mathbb H$ liegt.}
\end{thm}
\begin{proof}
\en{Choose}\de{Wähle den Faktor} $a=\frac{1}{\omega_1}$, \en{then we get}\de{dann gilt} $L'=a\cdot L = \mathbb Z +\mathbb Z \cdot \frac{\omega_2}{\omega_1}$.
\en{If}\de{Falls} $\im\lk\frac{\omega_2}{\omega_1}\rk>0$\en{, then we set}\de{ ist, setzen wir} $\tau=\frac{\omega_2}{\omega_1}$.
\en{If}\de{Falls} $\im\lk\frac{\omega_2}{\omega_1}\rk<0$\en{, then we set}\de{ ist, setzen wir} $\tau=-\frac{\omega_2}{\omega_1}$ \en{(this is still the same lattice, only another basic period).}\de{(das ist immer noch das gleiche Gitter, nur ein anderer Basisvektor).}
\en{The case}\de{Der Fall} $\im\lk\frac{\omega_2}{\omega_1}\rk=0$ \en{is impossible, since it would yield}\de{ist ausgeschlossen, weil sonst} $\frac{\omega_2}{\omega_1}\in\mathbb R$ \en{and $L$ wouldn't be a lattice (cf.~Def.~\ref{\en{EN}defgitter}).}\de{wäre und somit $L$ kein Gitter wäre (vgl.~Def.~\ref{\en{EN}defgitter}).}
\end{proof}

\begin{defi}\label{\en{EN}defimodtau}
\en{We call $\tau_1\in \mathbb H$ and $\tau_2\in \mathbb H$ "equivalent", iff the lattices $L_{\tau_1}$ and $L_{\tau_2}$ are equivalent.
For example, $\tau$ and $\tau+1$ are equivalent (since they generate the same lattice), but also $\tau$ and $-1/\tau$ are equivalent (since it holds $L_{-1/\tau}=1/\tau\cdot L_\tau$).
This explains why such equivalent $\tau_1$ and $\tau_2$ are also called "equivalent under modular transformations".}%
\de{Wir nennen $\tau_1\in \mathbb H$ und $\tau_2\in \mathbb H$ \glqq äquivalent\grqq, wenn die Gitter $L_{\tau_1}$ und $L_{\tau_2}$ äquivalent sind.
Beispielsweise sind $\tau$ und $\tau+1$ äquivalent (weil sie das gleiche Gitter erzeugen), aber auch $\tau$ und $-1/\tau$ sind äquivalent (weil $L_{-1/\tau}=1/\tau\cdot L_\tau$ ist).\\
Statt \glqqäquivalent\grqq~sagt man daher auch \glqqäquivalent unter Modultransformationen\grqq.}
\end{defi}

\begin{defi}\label{\en{EN}defijdelta}
\en{Given a lattice $L\subset\mathbb C$. Using the definitions of $g_2(L)$ and $g_3(L)$ from Prop.~\ref{\en{EN}dglP} we define the "discriminant" $\Delta$ and Klein's absolute invariant $J$ of the lattice:}%
\de{Gegeben sei ein Gitter $L\subset\mathbb C$. Mit den Definitionen von $g_2(L)$ und $g_3(L)$ aus Satz~\ref{\en{EN}dglP} definieren wir die \glqq Diskriminante\grqq\ $\Delta$ des Gitters und die \glqq absolute Invariante\grqq\ $J$ des Gitters:}
\begin{align*}
    \Delta(L) &:= g_2^3(L)-27g_3^2(L)\\
    J(L)&:=\frac{g_2^3(L)}{g_2^3(L)-27g_3^2(L)}
\end{align*}
\end{defi}

\begin{bem}
\en{If the lattice is of the form $L_\tau=\mathbb Z + \mathbb Z \tau$, we denote $g_2(\tau)$ instead of $g_2(L_\tau)$. In the same way, we write $g_3(\tau)$, $G_k(\tau)$, $\Delta(\tau)$ and $J(\tau)$.}%
\de{Wenn wir uns auf ein Gitter der Form $L_\tau=\mathbb Z + \mathbb Z \tau$ beziehen, schreiben wir kurz $g_2(\tau)$ statt $g_2(L_\tau)$. Ebenso schreiben wir abkürzend $g_3(\tau)$, $G_k(\tau)$, $\Delta(\tau)$ und $J(\tau)$.}
\end{bem}

\begin{thm}\label{\en{EN}trafog23}
\en{If the lattice $L'=a\cdot L$ is equivalent to $L$, then the following transformation formula for the Eisenstein series holds for $a\neq 0$:}%
\de{Wenn $L'=a\cdot L$ ein zu $L$ äquivalentes Gitter mit $a\neq 0$ ist, dann gilt folgende Transformationsformel für die Eisensteinreihen:}
$$G_k(aL) = a^{-k}\cdot G_k(L)$$
\en{and thus:}\de{und folglich:}
$$g_2(aL)=a^{-4}g_2(L)\qquad\text{\en{and}\de{und}}\qquad g_3(aL)=a^{-6}g_3(L)$$
\en{From this we get}\de{Hieraus wiederum folgt}
$$\Delta(aL)=a^{-12}\Delta(L)\qquad\text{\en{and}\de{und}}\qquad J(aL)=J(L)$$
\en{In particular Klein's absolute invariant $J$ has the same value if the lattices are equivalent -- this is why $J$ is called "invariant".}%
\de{Insbesondere ändert sich der Wert der absoluten Invariante $J$ nicht, wenn man zu einem äquivalenten Gitter übergeht. Das rechtfertigt den Namen Invariante.}
\end{thm}
\begin{proof}
\en{This is a consequence of the Def.~\ref{\en{EN}defeisen} of the Eisenstein series:}%
\de{Das ist eine Folge aus der Definition~\ref{\en{EN}defeisen} der Eisensteinreihen:}
$$G_k(aL)=\sum_{\substack{\omega'\in aL\\ \omega'\neq 0}}\omega'^{-k}=\sum_{\substack{\omega'\in aL\\ \omega'\neq 0}}\left(\frac{\omega'} a\right)^{-k}\cdot a^{-k} = a^{-k}\cdot \sum_{\substack{\omega\in L\\ \omega\neq 0}}\omega^{-k}= a^{-k}\cdot G_k(L)$$
\en{Here we used}\de{wobei} $\omega'=a\cdot\omega$\en{. This yields, with the definitions of $g_{2;3}$ from Prop.~\ref{\en{EN}dglP}, that}\de{ verwendet wurde. Es folgt} $g_2(aL)=60G_4(aL)=a^{-4}\cdot g_2(L)$\en{ and}\de{\linebreak
sowie} $g_3(aL)=140G_6(aL)= a^{-6}\cdot g_3(L)$.
\en{Finally, we get the discriminant}\de{Schließlich erhalten wir die Diskriminante}
$\Delta(aL)=(a^{-4})^3g_2^3(L) - 27 (a^{-6})^2g_3^2(L)=a^{-12}\cdot\Delta(L)$ \en{and Klein's absolute invariant $J(aL)=J(L)$, which shows that it doesn't change when the lattice is rotated and/or stretched.}\de{und die absolute Invariante $J(aL)=J(L)$, die sich also bei einer Drehstreckung des Gitters nicht ändert.}
\end{proof}

\begin{thm}\label{\en{EN}etatransf}
\en{For the basic periods and basic quasi periods of $L'=a\cdot L$ it holds:}%
\de{Für die Perioden und Quasiperioden von $L'=a\cdot L$ gilt:}
$$\omega_k'=a\cdot\omega_k\qquad\text{\en{and}\de{und}}\qquad \eta_k(L') = \frac 1 a \cdot \eta_k(L).$$
\end{thm}
\begin{proof}
\en{The first identity is proven by multiplicating the lattice with $a$.
Then, by Def.~\ref{\en{EN}defetak} and Prop.~\ref{\en{EN}eta}, it holds for any $z\in\mathbb C\setminus(L\cup L')$:}%
\de{Die erste Gleichung folgt aus der Multiplikation des Gitters mit der Zahl $a$.
Weiter folgt mit Def.~\ref{\en{EN}defetak} und Satz~\ref{\en{EN}eta} für beliebiges $z\in\mathbb C\setminus(L\cup L')$:}
\begin{align*}
    \eta_k(L') &= \eta_k(aL) = \zeta(z+a\omega_k;aL)-\zeta(z;aL) = \zeta(az+a\omega_k;aL)-\zeta(az;aL)
\end{align*}
\en{Next we use the Definition~\ref{\en{EN}defizeta} of the Weierstraß $\zeta$-function and get:}%
\de{Dann verwenden wir die Definition~\ref{\en{EN}defizeta} der Weierstraß'schen $\zeta$-Funktion und erhalten:}
\begin{align*}
\zeta(az;aL) &= \frac 1 {az} + \sum_{\substack{\omega\in aL\\ \omega\neq 0}} \left(\frac 1 {az-\omega} +\frac 1 \omega +\frac {az} {\omega^2}\right)
\end{align*}
\en{Now we change summation variables by setting $v := \omega/a$. Then, from $\omega\in aL$, we get $v\in L$ and thus:}%
\de{Jetzt folgt ein Variablenwechsel $v := \omega/a$. Aus $\omega\in aL$ folgt dann $v\in L$ und somit:}
\begin{align*}
    \zeta(az;aL)&=\frac 1 {az} + \sum_{\substack{v\in L\\ v\neq 0}} \left(\frac 1 {az-av} +\frac 1 {av} +\frac {az} {(av)^2}\right)=\frac 1 a \zeta(z;L)
\end{align*}
\en{In the same way (i.e.~setting $v := \omega/a$) we get $\zeta(az+a\omega_k;aL)=\frac 1 a \zeta(z+\omega_k;L)$ and}%
\de{Derselbe Variablenwechsel liefert analog $\zeta(az+a\omega_k;aL)=\frac 1 a \zeta(z+\omega_k;L)$ und}\belowdisplayskip=-12pt
\begin{align*}
    \eta_k(L') &= \zeta(az+a\omega_k;aL)-\zeta(az;aL) = \frac 1 a \zeta(z+\omega_k;L)-\frac 1 a \zeta(z;L) = \frac 1 a \cdot \eta_k(L)
\end{align*}
\end{proof}

\begin{defi}\label{\en{EN}definLJ}
\en{Given the lattice $L_\tau=\mathbb Z  + \mathbb Z \tau$, we define the equivalent lattice $L_J$ by}%
\de{Gegeben ist das Gitter $L_\tau=\mathbb Z  + \mathbb Z \tau$. Dann definieren wir ein zu $L_\tau$ äquivalentes Gitter $L_J$ durch}
\begin{align*}
    L_J &:= \mu(\tau) \cdot L_\tau\qquad\text{ \en{with}\de{mit} }\qquad\mu(\tau):=\sqrt{\frac{g_3(L_\tau)}{g_2(L_\tau)}}
\end{align*}
\en{From chapter~\ref{\en{EN}kappicardfuchs} onward, we will denote the basic periods of $L_J$ with $(\Omega_1,\Omega_2)$, and the corresponding basic quasi periods $\eta_k(L_J)$ will be called $(H_1,H_2)$.}%
\de{Die Basisperioden von $L_J$ bezeichnen wir im Folgenden mit $(\Omega_1,\Omega_2)$ und die zugehörigen Basisquasiperioden $\eta_k(L_J)$ bezeichnen wir mit $(H_1,H_2)$.}
\end{defi}

\begin{bem}\label{\en{EN}bemlj}
\en{It doesn't matter which branch of the square root is being chosen when calculating $\mu(\tau)$, because the negated basic periods generate the same lattice:}%
\de{Es ist egal, für welchen Zweig der Quadratwurzel man sich bei $\mu(\tau)$ entscheidet, denn negierte Basisperioden erzeugen das gleiche Gitter:}
$$\mathbb Z \omega_1 + \mathbb Z \omega_2 = \mathbb Z\cdot(-\omega_1) + \mathbb Z\cdot(-\omega_2).$$
\end{bem}

\begin{thm}\label{\en{EN}satz44}
\en{The plane affine algebraic curve $X(L_J)$ has a representation that depends only on the value of Klein's absolute invariant $J$ (that's why the lattice is called $L_J$). This representation reads:}%
\de{Die ebene affine Kurve zum Gitter $L_J$ hat eine Darstellung, die nur von der absoluten Invarianten $J$ des Gitters $L_J$ abhängt (deshalb nennt man das Gitter $L_J$). Diese lautet:}
$$X(L_J)=\left\{~(x,y)\in\mathbb C^2~\left|~ y^2 = 4 x^3 - \frac{27J}{J-1}( x +1)\right.~\right\}$$
\end{thm}
\begin{proof}
\en{From the transformation formula of $g_2$ and $g_3$ in Prop.~\ref{\en{EN}trafog23} and using $L_J=\mu(\tau)\cdot L_\tau$ we get:}%
\de{Aus dem Transformationsverhalten der $g_2$ und $g_3$ in Satz~\ref{\en{EN}trafog23} folgt mit $L_J=\mu(\tau)\cdot L_\tau$:}
\begin{align*}
        g_2(L_J) = \mu(\tau)^{-4}\cdot g_2(L_\tau) = \frac{g_2(L_\tau)^2}{g_3(L_\tau)^2}\cdot g_2(L_\tau) = \frac{g_2(L_\tau)^3}{g_3(L_\tau)^2}\\
        g_3(L_J) = \mu(\tau)^{-6}\cdot g_3(L_\tau) = \frac{g_2(L_\tau)^3}{g_3(L_\tau)^3}\cdot g_3(L_\tau) = \frac{g_2(L_\tau)^3}{g_3(L_\tau)^2}
\end{align*}
\en{In the lattice $L_J$ we thus have $g_2(L_J)=g_3(L_J)=:g$. Then we get the value of Klein's absolute invariant of $L_J$ from Def.~\ref{\en{EN}defijdelta}:}%
\de{Im Gitter $L_J$ gilt also $g_2(L_J)=g_3(L_J)=:g$. Die absolute Invariante des Gitters $L_J$ ist dann nach Definition~\ref{\en{EN}defijdelta}:}
$$J = \frac{g^3}{g^3-27g^2} = \frac{g}{g-27} \quad\Longrightarrow\quad g = \frac{27J}{J-1}$$
\en{which yields said equation of the plane affine algebraic curve $X(L_J)$:}%
\de{und wir erhalten die angekündigte ebene affine Kurve mit der Gleichung}\belowdisplayskip=-12pt
$$y^2 = 4 x^3 - g_2(L_J) x - g_3(L_J) = 4x^3 - g (x+1) = 4x^3 - \frac{27J}{J-1}(x+1).$$
\end{proof}
